\documentclass[12pt,a4paper]{article}

\usepackage[T1]{fontenc}
\usepackage[utf8]{inputenc}
\usepackage[margin=2.5cm]{geometry}
\usepackage{amsmath,amssymb,amsthm}
\usepackage{booktabs}
\usepackage[numbers,sort&compress]{natbib}
\usepackage[colorlinks=true,linkcolor=black,citecolor=black,urlcolor=black]{hyperref}
\usepackage{xcolor}
\usepackage{mdframed}
\usepackage{enumitem}
\usepackage{parskip}
\usepackage{setspace}
\usepackage{caption}

\newcommand{\doi}[1]{\href{https://doi.org/#1}{doi:#1}}

%% Reviewer comment box
\mdfdefinestyle{attackframe}{
  linecolor=red!60,
  linewidth=1pt,
  backgroundcolor=red!4,
  innertopmargin=10pt,
  innerbottommargin=10pt,
  innerleftmargin=14pt,
  innerrightmargin=14pt,
  skipabove=10pt,
  skipbelow=10pt
}

%% Defense box
\mdfdefinestyle{defenseframe}{
  linecolor=black!40,
  linewidth=0.8pt,
  backgroundcolor=black!3,
  innertopmargin=8pt,
  innerbottommargin=8pt,
  innerleftmargin=12pt,
  innerrightmargin=12pt,
  skipabove=6pt,
  skipbelow=6pt
}

\theoremstyle{definition}
\newtheorem{attack}{Attack}[section]

\title{\textbf{Adversarial Review: Kalaxi}\\\large
  Internal Red-Team Critique of\\
  \emph{``Kalaxi: A Constitutional Governance Architecture\\
  for Dignity-Constrained Software Systems''}\\[0.5em]
  \normalsize Phase V — Pre-Submission Stress Test}

\author{
  Prepared by: Voice (internal adversarial review)\\
  For: Mohamed Farag, KALAM\\
  \texttt{info@kalam.ch}
}

\date{February 2026}

\begin{document}

\maketitle
\thispagestyle{empty}

\begin{abstract}
This document is an internal adversarial review of the Kalaxi arXiv manuscript.
It is written from the perspective of a hostile but rigorous peer reviewer at a
top-tier venue (FAccT, ICML, NeurIPS, or IEEE S\&P). Its purpose is to surface every
weakness, inconsistency, overreach, and undefended claim before submission, so that
the paper can be strengthened or the claims properly bounded. Each attack is followed
by a recommended response — either a direct rebuttal, a paper revision, or an honest
acknowledgment. Nothing is minimized. The paper is strong. This document makes it
stronger.
\end{abstract}

\tableofcontents
\newpage

%% ═══════════════════════════════════════════════════════════════════════════
\section{Overview of Attack Surface}

The paper has five primary attack surfaces:

\begin{enumerate}
\item \textbf{Mathematical formalism:} Is the dignity predicate actually a
  well-defined mathematical object, or is it philosophy dressed in notation?
\item \textbf{Novelty claims:} Does Kalaxi contribute anything genuinely new,
  or is it a repackaging of existing governance frameworks?
\item \textbf{Empirical vacuity:} The paper has no experimental results.
  Is this a research paper or a design document?
\item \textbf{Practical viability:} Would any real institution deploy this?
  Is the architecture designed for the world that exists or a world that doesn't?
\item \textbf{Internal consistency:} Does the paper contradict itself?
  Are the theoretical claims coherent with the operational mechanisms?
\end{enumerate}

Each is attacked in turn. Severity ratings: \textbf{CRITICAL} (would cause rejection),
\textbf{MAJOR} (would require major revision), \textbf{MINOR} (would require minor
revision or clarification).

%% ═══════════════════════════════════════════════════════════════════════════
\section{Attack Surface 1: Mathematical Formalism}

\begin{attack}[The predicate is not a predicate — CRITICAL]
\label{att:notpredicate}
The paper defines $\mathcal{D}(e) = A(e) \times L(e) \times M(e)$ and calls
$A, L, M : \mathcal{E} \to \{0,1\}$ ``Boolean-valued functions.'' This notation
implies they are computable. But the component definitions reveal they are not.

``system\_reflects\_donor\_frame'' — how is this computed? What is a ``frame''?
How does a machine determine that it has reflected it? ``emotional\_signal\_recognized''
— this requires affect detection, a research problem with known unreliability across
demographic groups. ``no\_reduction\_to\_error\_object'' — how does a machine detect
that it has reduced a person to an error object? This is a social judgment requiring
rich contextual understanding that no current NLP system reliably provides.

\textbf{The paper claims a machine-evaluable predicate and then defines its components
in terms that are not machine-evaluable. This is not formalism — it is the appearance
of formalism.} The notation $A, L, M : \mathcal{E} \to \{0,1\}$ implies deterministic
computability. The components named do not have deterministic implementations.
\end{attack}

\begin{mdframed}[style=defenseframe]
\textbf{Recommended Response (paper revision required):}
This is the strongest attack and requires a direct response in the paper.
The paper must distinguish between:
(1) \emph{aspirational predicates} — the full philosophical definition;
(2) \emph{current operational approximations} — the heuristic implementations
in \texttt{dignity\_check.py};
(3) \emph{the gap between them} — which is itself an anomaly in the registry
(ANOM\#0004, GAP\#014, GAP\#015).

Add a subsection: ``Predicate Approximation and Its Limits.'' State clearly: the
Boolean formalization is a precision target, not a claim of current full
computability. Current implementations are heuristic approximations. The gap
between formalization and implementation is documented in the anomaly registry.
This is not a weakness to hide — it is the architecture's honest acknowledgment
of its frontier.

\textbf{Key sentence to add:} ``The Boolean formalization defines what the system
\emph{must} compute; current implementations approximate this via heuristics that
are explicitly logged as open gaps (GAP\#014, GAP\#015). The predicate's precision
is a research program, not a solved problem.''
\end{mdframed}

\begin{attack}[Non-compensability is trivially achievable by any veto — MAJOR]
\label{att:veto}
The paper presents non-compensability as a deep theoretical property motivated by
deontological ethics. But any AND-gate has this property. Connecting it to Kant is
motivated reasoning. A reviewer will note: ``The authors' formulation is equivalent
to requiring that all three Boolean conditions be TRUE before proceeding. This is
a logical AND. Every checklist system is 'non-compensatory' by this definition.
The Kantian framing adds philosophical weight but no structural novelty.''
\end{attack}

\begin{mdframed}[style=defenseframe]
\textbf{Recommended Response (clarification required):}
The reviewer is technically correct: the multiplicative form is equivalent to a
logical AND over Boolean variables. The paper must be explicit about this and
defend the choice on different grounds:

The non-compensability claim is not about the Boolean structure per se, but about
the \emph{choice of variables and the refusal to weight them}. Most governance
systems implicitly weight factors (A matters more than M in urgent cases, etc.).
Kalaxi refuses this weighting — all three conditions are equal in weight and any
one zero halts the system. \emph{That} is the deontological claim: not that the
logic is AND, but that the weights are equal and non-negotiable.

Also strengthen: the paper should note that the equivalent claim in weighted-sum
formulations would allow $A = 0, L = 1, M = 1, w_A = 0.2$ to produce a positive
dignity score. Kalaxi's structure explicitly forbids this class of override.
\end{mdframed}

\begin{attack}[The damping model is a metaphor, not a model — MAJOR]
\label{att:dampingmetaphor}
Section V presents a second-order differential equation for normative state evolution.
The paper gives no method for measuring $\omega_0$, $\zeta$, or $N(t)$ empirically.
There is no operational definition of ``normative state.'' There is no evidence that
governance dynamics actually follow second-order linear differential equations. The
use of Eq.~(3) is scientism: the equation gives the appearance of rigor without
delivering it.

A reviewer in a control-theory track will write: ``The authors assert that their
thermal delay implements critical damping ($\zeta \geq 1$) but provide no empirical
measurement of the system's actual damping ratio, no method for measuring normative
state $N(t)$, and no validation of the second-order linear model as an accurate
description of governance dynamics. The analogy to control systems is suggestive
but not substantiated.''
\end{attack}

\begin{mdframed}[style=defenseframe]
\textbf{Recommended Response (scope reduction required):}
This attack is valid. The differential equation should either be substantiated
or reframed. Recommended approach: reframe the control-theoretic analysis as
an \emph{analogy} that motivates the design, not a formal model of it.

Replace ``We model normative state evolution as a second-order system'' with
``We draw on the intuition of second-order damping to motivate the thermal delay
mechanism.'' Remove the differential equation from the main text; move it to an
appendix clearly labeled ``Motivating Analogy.'' State explicitly: ``We do not
claim that governance dynamics follow this differential equation, nor do we
provide empirical measurements of $\omega_0$ or $\zeta$. The second-order model
is a conceptual scaffold, not a validated empirical model.''

This is not a defeat — it is precision. The thermal delay is justified by
other arguments (the Zeno effect, the empirical observation of governance overshoot)
that do not require the differential equation.
\end{mdframed}

\begin{attack}[The collective dignity formula is dimensionally inconsistent — MINOR]
\label{att:collective}
Equation~(2): $\mathcal{D}_\text{cohort} = \bar{D}(\mathcal{C},T) \times
(1 - \sigma^2_\text{penalty}(\mathcal{C},T))$.

If $\sigma^2_\text{penalty} > 1$, the collective dignity becomes negative, which
is undefined. The paper does not specify the range of $\sigma^2_\text{penalty}$
or how it is computed from the variance of $\{D(e)\}$. Since $D(e) \in \{0,1\}$,
the variance of this Bernoulli distribution is $\bar{D}(1-\bar{D}) \leq 0.25$.
If the penalty is simply the variance, the formula is well-defined. If it is a
normalized or scaled version, this must be stated.
\end{attack}

\begin{mdframed}[style=defenseframe]
\textbf{Recommended Response (minor revision):}
Specify the computation of $\sigma^2_\text{penalty}$ precisely. Recommended:
define it as the normalized variance: $\sigma^2_\text{penalty} = 4 \cdot
\text{Var}(\{D(e)\})$, which has range $[0,1]$ when $D(e) \in \{0,1\}$.
This ensures $\mathcal{D}_\text{cohort} \in [0,1]$ and is well-defined.
Add a single sentence to the paper specifying the normalization.
\end{mdframed}

%% ═══════════════════════════════════════════════════════════════════════════
\section{Attack Surface 2: Novelty Claims}

\begin{attack}[This is just Value Sensitive Design with a Git wrapper — CRITICAL]
\label{att:vsd}
The paper's central claim is that it embeds governance into repository structure.
But Value Sensitive Design \cite{friedman2006} already proposes embedding values
into technical design. The NIST AI RMF \cite{nist2023} proposes governance
throughout the AI lifecycle. Constitutional AI \cite{bai2022} embeds principles
as training objectives. The EU AI Act \cite{euaiact2021} mandates governance
documentation.

A reviewer will ask: ``What does Kalaxi add that these frameworks do not already
provide? The authors' answer — that Kalaxi makes governance `executable' via
version control — is not clearly distinguished from existing approaches. GitHub
Actions can enforce policy checks. Existing CI/CD pipelines already block merges
on policy violations. What is genuinely new here?''
\end{attack}

\begin{mdframed}[style=defenseframe]
\textbf{Recommended Response (paper revision required):}
This is the second strongest attack. The paper needs a sharper novelty statement.
The genuine novelty claims are:

\begin{enumerate}
\item \textbf{The predicate is pre-output, not post-hoc.} VSD embeds values in
  design process. NIST RMF embeds governance in documentation. Constitutional AI
  embeds principles in training. Kalaxi's dignity gate fires \emph{before every
  output is rendered} — it is an operational block, not a design commitment or a
  training objective.

\item \textbf{The thermal delay is an ethical primitive.} No existing governance
  framework treats temporal delay as a first-class ethical mechanism. The Zeno
  protection justification is novel.

\item \textbf{The append-only proverb canon as governance artifact.} The
  paremiologically-grounded wisdom layer — 3,333 proverbs linked to 1,100+
  structured anomalies via a ratified UCS graph — has no precedent in governance
  literature. It is not documentation. It is a growing normative knowledge base
  with formal integrity constraints (COV\#006: every anomaly linked to a proverb).

\item \textbf{The steward's mirror ritual as governance mechanism.} Self-reflexive
  logging as a structural requirement for ratification authority is not present in
  existing frameworks.
\end{enumerate}

Add a ``Novelty Clarification'' subsection in Related Work or Contributions that
directly names these four distinctions and explains why each represents a structural
advance over prior work.
\end{mdframed}

\begin{attack}[The Quantum Zeno analogy is not a contribution — MINOR]
\label{att:zeno}
The paper introduces the Quantum Zeno effect as a novel justification for thermal
delay. A reviewer will note: ``The Quantum Zeno effect describes quantum measurement.
Applying it to governance is an analogy, not a theoretical result. The authors
provide no formal connection between quantum measurement collapse and normative
seed maturation. This is an evocative metaphor that does not strengthen the
technical argument.''
\end{attack}

\begin{mdframed}[style=defenseframe]
\textbf{Recommended Response (reframe as motivation):}
Agree with the reviewer. In the final paper, reframe: ``The quantum Zeno effect
provides an \emph{intuition} for why frequent observation can prevent development
of a seed's full semantic potential, analogous to how measurement collapses quantum
superposition.'' Do not claim this as a formal result. It is motivating intuition.
The primary justification for thermal delay remains the control-theoretic damping
argument and the empirical observation of governance overshoot — both stronger
footing.
\end{mdframed}

%% ═══════════════════════════════════════════════════════════════════════════
\section{Attack Surface 3: Empirical Vacuity}

\begin{attack}[There are no results — CRITICAL]
\label{att:noresults}
This is the most common and most damaging critique a systems/governance paper can
receive. The paper proposes an architecture, formally describes it, and then
describes an evaluation strategy. It does not report any evaluation results.

At FAccT, NeurIPS, or ICML, papers without results are rejected. At IEEE venues,
a system paper requires at minimum a feasibility demonstration. A reviewer will
write: ``The authors describe an interesting architecture but provide no evidence
that it functions as described. No case studies. No user studies. No simulations.
No measurement of whether the dignity predicate actually blocks harmful outputs.
No measurement of thermal delay's effect on governance stability. This is a design
proposal, not a research paper.''
\end{attack}

\begin{mdframed}[style=defenseframe]
\textbf{Recommended Response (venue selection and framing):}
This attack is decisive for venue selection. Three options:

\textbf{Option A — Submit to a systems/design track} where the contribution
is the architecture itself. Several venues accept this:
\emph{CHI} (design contributions), \emph{CSCW} (socio-technical systems),
\emph{IEEE Software} (architecture papers), \emph{SSRN} or \emph{arXiv}
preprint as position paper.

\textbf{Option B — Add a feasibility case study.} Use the hiring algorithm
simulation already described in the paper. Run it. The THRESHOLD.md already
contains ANOM\#PROXY-001, ANOM\#PROXY-CASCADE-001, and W\#HIRING-001. Deploy
the dignity predicate on a labeled hiring dataset. Report REFUSAL\_STATE
activation rates, cohort disparity before and after. Even a small simulation
provides the ``results'' section the paper currently lacks.

\textbf{Option C — Reframe as position/vision paper.} Some venues (FAccT,
AIES) accept position papers that argue for a new approach without full
experimental validation. Frame explicitly: ``This paper presents a new
\emph{paradigm} for governance as infrastructure. We report implementation
status and propose an evaluation agenda. Full empirical validation is future work.''
This is the honest frame and is accepted at some venues.

\textbf{Recommendation: Option B + C.} Add the hiring simulation (it is already
half-designed in the anomaly registry) and frame as a position/vision paper with
a proof-of-concept case study.
\end{mdframed}

\begin{attack}[The KPIs are not measurable as stated — MAJOR]
\label{att:kpis}
The paper lists KPIs including ``fraction of steward tending sessions recording
a `pause' or `smile' response (target $\geq 30\%$).'' A `pause' is not a
measurable unit. A `smile' requires a camera and facial recognition. These are
not scientific KPIs. Similarly, ``steward surprise rate'' is entirely subjective.

The paper's KPIs mix objective metrics (REFUSAL\_STATE activation rate, cohort
D variance) with phenomenological observations (pause, smile) without
distinguishing them or explaining how the latter are operationalized.
\end{attack}

\begin{mdframed}[style=defenseframe]
\textbf{Recommended Response (revision required):}
Separate KPIs into two tiers:

\textbf{Tier A — Objective, machine-measurable:} REFUSAL\_STATE activation
rate; cohort D variance ($\sigma^2_\text{cohort}$); donor flag rate on
humour seeds; thermal delay violation rate (seeds ratified before delay
expired); override logging compliance rate.

\textbf{Tier B — Qualitative, steward-assessed:} Proverb resonance
(steward-reported, using a structured 5-point Likert scale rather than
``pause or smile''); cross-anomaly compression quality (evaluated by
steward against Lock Test criteria on a structured rubric).

Remove the language about ``smile'' and ``pause'' from quantitative KPI
lists. Retain them in the qualitative methodology section with explicit
acknowledgment that they are phenomenological observations requiring a
structured observation protocol, not automated metrics.
\end{mdframed}

%% ═══════════════════════════════════════════════════════════════════════════
\section{Attack Surface 4: Practical Viability}

\begin{attack}[No real institution would deploy this — MAJOR]
\label{att:noinstitution}
The architecture requires: a single dedicated human (the steward) who performs
a ritual reflection before every governance action; a 7--90 day delay before any
normative update; a human who reads and ratifies every new proverb; an
append-only repository that grows indefinitely; and a cultural foundation
(Hakaka narrative, Swiss institutional context) that may not translate.

A reviewer in governance or public policy will note: ``The steward model is
incompatible with institutional realities. Healthcare systems, credit bureaus,
and employment platforms process millions of decisions daily. A 7-day delay
before any normative update, and a single human ratification authority, would
paralyze operations. The paper addresses scalability as a limitation but does
not provide a credible path to addressing it.''
\end{attack}

\begin{mdframed}[style=defenseframe]
\textbf{Recommended Response (scope clarification required):}
The paper needs to be explicit about deployment scope. Kalaxi is not designed
for high-frequency automated decision systems processing millions of transactions
per minute. It is designed for:

\begin{enumerate}
\item \textbf{High-stakes, lower-frequency decisions} — hiring, benefit
  allocation, educational assessment — where a 7-day normative update delay
  is operationally acceptable.
\item \textbf{Governance of the governance layer} — not the individual decisions
  themselves but the normative framework within which those decisions are made.
  The dignity predicate runs pre-output; the thermal delay governs updates to
  the predicate, not the predicate's execution.
\item \textbf{Institutional pilots} — Kalaxi's current deployment context is
  educational and community governance settings in Switzerland, where the
  steward model is feasible.
\end{enumerate}

Add to the paper: ``Kalaxi is not designed for high-frequency transactional
systems. It targets governance of the normative framework within which
decisions are made, not real-time evaluation of each decision. The dignity
predicate runs pre-output; thermal delay governs updates to the predicate
itself. In practice, normative updates to governance frameworks occur
monthly or quarterly, making 7-day delays operationally acceptable in
most target contexts.''
\end{mdframed}

\begin{attack}[The steward is a single point of failure with no true mitigation — MAJOR]
\label{att:steward}
The paper acknowledges the steward as a single point of failure and proposes a
Council Ratification Extension as mitigation — but this is described as ``under
development.'' The Mirror Ritual is a social obligation, not a technical
enforcement. The Override Logging requires the steward to log their own
violations, which is the governance equivalent of self-policing.

A reviewer will note: ``The paper's principal-agent analysis correctly identifies
moral hazard in the steward relationship but then proposes mitigations that are
themselves vulnerable to the same hazard. A steward who chooses not to perform
the Mirror Ritual can simply not perform it. A steward who makes an override can
choose not to log it. The paper does not explain how these mitigations have
technical enforcement.''
\end{attack}

\begin{mdframed}[style=defenseframe]
\textbf{Recommended Response (revision required):}
This attack is partially valid. The paper should distinguish more carefully between
\emph{social enforcement} and \emph{technical enforcement}:

\textbf{Technically enforced:} The CI/CD gate described in Section VIII rejects
commits to the canonical branch unless \texttt{mirror.md} has a timestamp within
24 hours. This is machine-enforced, not social. The append-only audit means that
\emph{if} the steward commits without a mirror entry, this absence is detectable
in the audit trail. It is not invisible.

\textbf{Socially enforced:} The \emph{quality} of the mirror entry is social.
A steward who writes ``I am fine'' every day without genuine reflection satisfies
the machine check but not the spirit.

The honest position: ``Technical enforcement ensures the Mirror Ritual \emph{occurs};
social accountability ensures it \emph{means something}. This is analogous to
mandatory disclosure requirements in financial regulation — the requirement to
disclose is technically enforced; the quality of disclosure is socially audited.''

Add this distinction to the paper. It strengthens rather than weakens the argument.
\end{mdframed}

\begin{attack}[The cultural signature is a deployment barrier — MINOR]
\label{att:cultural}
The narrative foundation (Hakaka, Swiss context, particular proverb tradition)
creates a cultural specificity that may limit adoption. A reviewer in a global
or post-colonial computing context will flag: ``The system's legitimacy claim
depends on donor recognition — people seeing themselves in the system. But the
system's founding narrative, proverb tradition, and institutional context are
culturally particular. A Kenyan healthcare system or a Brazilian educational
platform may not find their experience reflected in a Swiss-contextured narrative
about knot-tying ancestors. The claim to universality through the UDHR mapping
requires stronger cross-cultural validation.''
\end{attack}

\begin{mdframed}[style=defenseframe]
\textbf{Recommended Response (acknowledgment):}
This is a fair and important point. Do not minimize it. Add to the epistemic
limitations section: ``The cultural signature of Kalaxi's founding narrative
is a genuine barrier to universal adoption. We do not claim the Hakaka narrative
is universal; we claim the \emph{structural pattern} it encodes — survival under
constraint, dignity as the first requirement, anomaly as the system's most honest
signal — is universal. The cross-cultural proverb recognition study proposed in
Section~XII is designed specifically to test whether this structural pattern
translates across cultural contexts. Until that data exists, the universality
claim is theoretical.''
\end{mdframed}

%% ═══════════════════════════════════════════════════════════════════════════
\section{Attack Surface 5: Internal Consistency}

\begin{attack}[The paper claims both deontological and consequentialist foundations — contradiction or confusion? — MAJOR]
\label{att:ethics}
Section VII argues that Kalaxi ``deliberately occupies both'' deontological and
consequentialist traditions. A reviewer in philosophy of ethics will challenge:
``Deontological and consequentialist ethics are not merely different emphases —
they are systematically incompatible in cases of genuine conflict. If the sealed
gate prohibits an action that would maximize aggregate welfare (the classic
trolley-problem structure), the architecture is deontological, full stop. The
claim to occupy both traditions simultaneously is incoherent unless the paper
specifies exactly how conflicts between them are resolved.''
\end{attack}

\begin{mdframed}[style=defenseframe]
\textbf{Recommended Response (clarification required):}
The paper needs to be precise about where the traditions apply and how conflicts
are resolved:

The sealed gate applies \emph{when the deontological constraint is triggered}.
In that case, the constraint is absolute and no consequentialist override exists.
The system is deontological in those cases.

Within the space of actions that do not trigger the sealed gate, the system
optimizes consequentially — it grows its wisdom canon toward better outcomes.

The synthesis is not that both traditions apply simultaneously to the same
decision. It is that different tiers of the architecture apply different ethical
frameworks: Tier-1 (Stone) is deontological; within that floor, the system is
consequentialist. This is Rawls's original structure: the priority of the right
over the good, but within that priority, optimization toward the good.

Revise Section VII to state this priority ordering explicitly: ``When sealed-gate
conditions are triggered, the architecture is deontological and admits no
consequentialist override. Within the space of actions not triggering the sealed
gate, the architecture optimizes consequentially. The two traditions do not
compete; they apply at different tiers.''
\end{mdframed}

\begin{attack}[The proverb canon is claimed as governance infrastructure but is not technically load-bearing — MAJOR]
\label{att:provinfra}
The paper claims the 3,333 proverbs are ``load-bearing infrastructure rather
than decoration'' (citing COV\#006). But COV\#006 merely requires that every
anomaly be \emph{linked to} a proverb. It does not specify how proverbs are
\emph{used} in governance decisions. No module description shows proverbs
being consulted at decision time. No mechanism routes a proverb to an output
pathway. The proverb-anomaly link is a metadata annotation, not an operational
dependency.

A reviewer will write: ``The paper claims proverbs are operational infrastructure.
But the module descriptions show that the dignity predicate, the Check module,
and the Say module operate on Boolean conditions and covenant checks, not on
proverbs. Proverbs appear to be documentation.''
\end{attack}

\begin{mdframed}[style=defenseframe]
\textbf{Recommended Response (revision or honest reframing):}
This is a valid structural critique. Two options:

\textbf{Option A — Strengthen the operational claim.} Add a mechanism by which
proverbs are operationally consulted. The Weave module can be specified to use
the proverb canon as a target space for cross-anomaly compression: a candidate
proverb is evaluated against the existing canon for redundancy and resonance
before being proposed. This makes the canon load-bearing in pattern synthesis.

\textbf{Option B — Reframe honestly.} If proverbs are currently metadata
annotations with a mandate for coverage (COV\#006) but not operational
decision-making roles, say so. ``The proverb canon is \emph{normatively}
load-bearing: every anomaly must be linked to a proverb that names its
structural pattern, making the system's anomaly registry interpretable
by humans rather than only by machines. The operational load-bearing role
of proverbs in real-time decision-making is future work.''

Recommended: Option A for the cross-anomaly compression mechanism
(already specified), plus Option B's honest reframe for real-time
decision-making.
\end{mdframed}

\begin{attack}[The append-only claim is not technically substantiated — MINOR]
\label{att:appendonly}
The paper claims append-only audit with Merkle-chained receipts. Git does not
natively produce Merkle-chained receipts in the format described. Git's commit
graph is a DAG with SHA-1/SHA-256 hashes, but the paper's description of
``Merkle-chained receipts'' implies a specific data structure that is not described
technically. A reviewer with systems security background will ask: ``What is
the exact data structure of the Merkle-chained receipt? How are collisions
prevented? What is the threat model for an attacker who has write access to
the repository?''
\end{attack}

\begin{mdframed}[style=defenseframe]
\textbf{Recommended Response (precision required):}
Either (a) specify the Merkle chain structure in an appendix, or (b) replace
``Merkle-chained'' with the accurate description: ``Git's native SHA-256 commit
DAG, where each commit contains the hash of its parent, making the full history
tamper-evident under standard Git security assumptions.''

Also add: ``Under the assumption that the repository host (GitHub or equivalent)
is trusted, the commit history is tamper-evident. Against an attacker with
repository write access and the ability to force-push, the append-only guarantee
is social rather than cryptographic.'' This is the honest scope of the security
claim.
\end{mdframed}

%% ═══════════════════════════════════════════════════════════════════════════
\section{Summary: Priority Revision Checklist}

\subsection*{Before Submission — Must Fix}

\begin{enumerate}[label=\textbf{M\arabic*.},leftmargin=*]
\item \textbf{Predicate computability.} Add ``Predicate Approximation and Its
  Limits'' subsection. Distinguish aspiration from current implementation.
  Cite GAP\#014 and GAP\#015 as the canonical self-criticism.
  (Addresses Attack~\ref{att:notpredicate})

\item \textbf{Novelty statement.} Add dedicated ``Novelty Clarification''
  subsection in Related Work. State four distinct contributions vs.\ prior
  work. Direct answer to VSD/NIST/Constitutional AI comparisons.
  (Addresses Attack~\ref{att:vsd})

\item \textbf{Empirical content.} Either add hiring simulation results or
  reframe explicitly as position/vision paper with proof-of-concept.
  Choose venue accordingly.
  (Addresses Attack~\ref{att:noresults})

\item \textbf{Damping model reframe.} Move differential equation to appendix
  labeled ``Motivating Analogy.'' State explicitly: not a validated empirical
  model. The thermal delay is justified by control-theoretic intuition and
  the Zeno analogy, not by a fitted second-order system.
  (Addresses Attack~\ref{att:dampingmetaphor})

\item \textbf{Ethical framework priority ordering.} Revise Section VII to
  state clearly: Tier-1 is deontological (no consequentialist override);
  within Tier-1, the system is consequentialist. The traditions apply at
  different tiers, not simultaneously to the same decision.
  (Addresses Attack~\ref{att:ethics})

\item \textbf{Deployment scope clarification.} Add explicit paragraph stating
  that Kalaxi targets high-stakes, lower-frequency governance contexts, not
  high-throughput transactional systems. Thermal delay applies to normative
  updates, not to predicate execution.
  (Addresses Attack~\ref{att:noinstitution})
\end{enumerate}

\subsection*{Should Fix — Strengthens the Paper}

\begin{enumerate}[label=\textbf{S\arabic*.},leftmargin=*]
\item \textbf{Proverb operational role.} Specify how the proverb canon is
  operationally used in cross-anomaly compression. Add mechanism description.
  (Addresses Attack~\ref{att:provinfra})

\item \textbf{Steward enforcement distinction.} Distinguish technical
  enforcement (CI gate) from social enforcement (quality of mirror entry).
  Use financial disclosure analogy.
  (Addresses Attack~\ref{att:steward})

\item \textbf{Collective dignity formula.} Specify $\sigma^2_\text{penalty}$
  normalization. One sentence.
  (Addresses Attack~\ref{att:collective})

\item \textbf{Non-compensability clarification.} Explain why AND-gate + equal
  weights is deontologically meaningful, not just an AND-gate.
  (Addresses Attack~\ref{att:veto})

\item \textbf{KPI tier separation.} Separate objective metrics from
  qualitative observations. Remove ``smile'' from quantitative KPI list.
  (Addresses Attack~\ref{att:kpis})

\item \textbf{Append-only security scope.} Specify Git SHA-256 DAG as the
  actual mechanism. State the trust assumption (trusted repository host).
  (Addresses Attack~\ref{att:appendonly})
\end{enumerate}

\subsection*{Acknowledge and Leave — Honest Limitations}

\begin{enumerate}[label=\textbf{A\arabic*.},leftmargin=*]
\item \textbf{Quantum Zeno reframe.} Present as motivating intuition,
  not formal result.
  (Addresses Attack~\ref{att:zeno})

\item \textbf{Cultural signature.} Name explicitly in limitations. Propose
  cross-cultural validation study. Do not minimize.
  (Addresses Attack~\ref{att:cultural})
\end{enumerate}

%% ═══════════════════════════════════════════════════════════════════════════
\section{What the Paper Does Well — Retain These}

The adversarial review is complete. In fairness, here is what the paper does
genuinely well and should not be changed under reviewer pressure:

\begin{enumerate}
\item \textbf{The epistemic limitations section.} It is unusual and valuable.
  Most papers hide limitations; this one names them with precision and
  connects them to the anomaly registry. Reviewers who notice this will
  respect it.

\item \textbf{The non-compensatory predicate as structural claim.} Despite
  Attack~\ref{att:veto}'s correct observation, the equal-weight no-override
  structure is a genuine policy commitment that most governance frameworks
  do not make explicitly. It is worth defending.

\item \textbf{The narrative foundation.} It is unusual. Some reviewers will
  be hostile; others will find it memorable. It is part of the paper's
  identity. The Acknowledgments section (``To Laila, Yara, and Salim'') is
  one of the most humanizing acknowledgments in an AI governance paper and
  should stay.

\item \textbf{The formal properties appendix.} Properties 1--4 of the
  dignity predicate are precisely stated and correct. The determinism and
  transparency properties in particular are stronger claims than most
  governance papers make.

\item \textbf{The principal-agent analysis.} The two-level hierarchy
  (donors $\to$ steward $\to$ modules) is clearly stated and the moral
  hazard mitigations are concrete. This section is the paper's best
  theoretical contribution.
\end{enumerate}

%% ═══════════════════════════════════════════════════════════════════════════
\section{Venue Recommendations}

Based on this analysis:

\begin{description}
\item[\textbf{Optimal venue (no added results):}] \emph{FAccT 2027} (ACM Conference
  on Fairness, Accountability, and Transparency) — Vision/Position paper track.
  Or \emph{AIES 2026} (AAAI/ACM Conference on AI, Ethics, and Society). Both
  accept system design papers with honest empirical limitations.

\item[\textbf{Optimal venue (with hiring simulation):}] \emph{FAccT 2026}
  full paper track, or \emph{CSCW 2026} (Computer-Supported Cooperative Work
  and Social Computing). With a case study, the empirical gap closes enough
  for a full paper submission.

\item[\textbf{arXiv preprint (now):}] Submit the Phase IV version to
  cs.AI and cs.CY (Computers and Society) immediately. Establish priority.
  The arXiv version is submission-ready. Peer review strengthens it;
  priority requires it now.

\item[\textbf{Avoid:}] NeurIPS, ICML, ICLR (require strong empirical
  results in ML). IEEE S\&P (requires security proof). \emph{Nature Machine
  Intelligence} (requires large-scale deployment evidence).
\end{description}

%% ═══════════════════════════════════════════════════════════════════════════
\section{Final Assessment}

\begin{mdframed}[style=defenseframe]
\textbf{Honest summary:}

The paper makes a genuine contribution to AI governance architecture. The
principal-agent analysis, the non-compensatory predicate, the thermal delay
mechanism, and the anomaly-proverb linkage are structurally sound and well-motivated.

The paper's primary weaknesses are: (1) the gap between formal notation and
current computational implementation; (2) the absence of empirical results;
(3) insufficient clarity on deployment scope.

All three are addressable without changing the paper's core claims.
None of them invalidate the architecture.

\textbf{Probability of acceptance at FAccT Vision track (as-is):} $\approx$40\%.\\
\textbf{Probability of acceptance after M1--M6 revisions:} $\approx$70\%.\\
\textbf{Probability of acceptance at FAccT full paper (with hiring simulation):}
$\approx$55--65\%.

The paper should be submitted to arXiv now. The revisions are a week of work.
The architecture is ready.
\end{mdframed}

\bigskip
\noindent\textit{The stone that knows its cracks does not break — it knows
where to hold.}

%% ─────────────────────────────────────────────────────────────────────────
\bibliographystyle{plainnat}
\begin{thebibliography}{10}

\bibitem{friedman2006}
B.~Friedman, P.~H. Kahn~Jr., and A.~Borning,
``Value sensitive design and information systems,''
in \emph{Human-Computer Interaction and Management Information Systems:
Foundations}, 2006.

\bibitem{bai2022}
Y.~Bai et al.,
``Constitutional AI: Harmlessness from AI feedback,''
\emph{arXiv preprint arXiv:2212.08073}, 2022.

\bibitem{nist2023}
NIST,
``Artificial Intelligence Risk Management Framework (AI RMF 1.0),''
NIST AI 100-1, 2023.

\bibitem{euaiact2021}
European Commission,
``Proposal for a Regulation: Artificial Intelligence Act,''
COM(2021) 206 final, 2021.

\end{thebibliography}

\end{document}
